\chapter{Introduction}
\label{ch:introduction}

\chapterepigraph{The sea, once it casts its spell, holds one in its net of
wonder forever.}{Jacques-Yves Cousteau}

\section{Blue carbon and the question of permanence}
\label{sec:permanence}

Vegetated coastal ecosystems occupy less than half a per cent of the sea
floor, yet they account for a disproportionate share of the organic
carbon buried in marine sediments \parencite{rowan2017seagrass}. Seagrass
meadows are the most extensive of these systems in temperate waters. Their
canopies slow the water, trap suspended particles and bury both their own
detritus and material imported from land; low oxygen beneath the rhizome
mat then slows decomposition to a crawl. The result is a sediment store
that, in sheltered Scandinavian estuaries, has accumulated without
interruption for at least two centuries \parencite{holmgren2019burial}.

This store has become the basis of an emerging market in \enquote{blue
carbon} credits, in which the restoration or protection of a meadow is
credited with the carbon it buries. National inventory guidelines treat
that carbon as permanently removed once it lies below the surface mixed
layer \parencite{bcwg2023guidelines}. The assumption is convenient but
largely untested. A meadow that loses its canopy also loses the
mechanism that protects its sediment, and a recent review found that
disturbance can return between a tenth and a half of the upper sediment
carbon to the water column within a few years
\parencite{fairbairn2022permanence}.

The long-term balance of a meadow can be written as a simple budget. If
\(B\) is the burial flux, \(R\) the remineralisation flux and \(E\) the
loss by erosion, the change in the sediment carbon stock \(C\) per unit
area is
\begin{equation}
  \frac{\mathrm{d}C}{\mathrm{d}t} = B(t) - R(t, T) - E(t),
  \label{eq:budget}
\end{equation}
where the explicit dependence of \(R\) on water temperature \(T\) is the
subject of this thesis. Under a stable canopy \(B\) dominates and the
store grows; the question is what happens to the sign of
\cref{eq:budget} when the canopy collapses during a heatwave.

\section{Marine heatwaves in shallow seas}
\label{sec:heatwaves}

Marine heatwaves, periods of at least five days in which sea temperature
exceeds the ninetieth percentile of its seasonal climatology, have become
both longer and more frequent across the shallow seas of northern Europe
\parencite{marsh2021heatwaves}. Shallow estuaries are particularly
exposed: their small volume warms quickly under calm, sunny weather, and
the eelgrass that grows there already lives close to its upper thermal
limit in late summer \parencite{okafor2018thermal}.

\begin{figure}[tb]
  \centering
  \begin{tikzpicture}[x=1mm, y=1mm, font=\footnotesize,
      flux/.style={-{Stealth[length=2mm]}, line width=0.6pt}]
    % water and sediment
    \fill[accent!6] (0,0) rectangle (100,30);
    \fill[black!8] (0,-22) rectangle (100,0);
    \draw[hairline, line width=0.4pt] (0,0) -- (100,0);
    \draw[hairline, decorate, decoration={snake, amplitude=0.6mm,
      segment length=8mm}] (0,30) -- (100,30);
    % shoots
    \foreach \x in {14,20,...,62} {
      \draw[accent!70!black, line width=0.7pt]
        (\x,0) .. controls (\x-1,8) and (\x+2,14) .. (\x+1,20);
    }
    \draw[black!45, line width=0.5pt] (10,-3) -- (66,-3);
    % fluxes
    \draw[flux] (82,26) -- (82,2) node[pos=0.1, right] {deposition};
    \draw[flux, accent] (38,-6) -- (38,-18) node[midway, right] {burial \(B\)};
    \draw[flux] (76,-14) -- (76,-2) node[midway, right] {remineralisation \(R\)};
    \draw[flux] (6,2) -- (6,12) node[midway, right] {erosion \(E\)};
    \node[anchor=west] at (68,22) {canopy};
    \node[anchor=west, black!60] at (2,-19) {sediment};
  \end{tikzpicture}
  \caption{The carbon budget of a seagrass meadow. The canopy traps
    particles and the rhizome mat protects buried carbon; both
    mechanisms weaken when shoot density falls.}
  \label{fig:budget}
\end{figure}

\Cref{fig:budget} summarises the mechanisms involved. When a heatwave
thins the canopy, deposition falls, the sediment surface is exposed to
waves and the oxygenated layer deepens, so all three terms of
\cref{eq:budget} move against the store at once. Whether this happens
gradually, in proportion to the heat received, or abruptly once a
threshold is crossed matters a great deal for management: a threshold
offers a target that can be monitored, while a gradual response does
not.

\section{Aims and structure of the thesis}
\label{sec:aims}

The thesis pursues three aims:
\begin{enumerate}
  \item to quantify how eelgrass shoot density responds to the cumulative
        heat dose of a summer, and to test for a threshold in that
        response;
  \item to measure the loss of sediment organic carbon that follows
        canopy decline, and its depth distribution;
  \item to estimate the consequences for blue-carbon accounting in the
        estuary as a whole.
\end{enumerate}
\Cref{ch:study-system} introduces the Westhaven estuary, the twelve
monitored meadows and the methods common to all chapters.
\Cref{ch:heatwave-response} reports the shoot and sediment responses and
the threshold model. Field protocols are collected in
\cref{app:protocols}.
